% GENERATED FILE. Edit canonical lessons, then run npm run generate:books.
% canonical-source-hash: fnv1a64:9fc5a57b71aed24f
% canonical-lessons: RU-C74-always, RU-C74-sometimes, RU-C74-never, RU-C74-now, RU-C74-later

\chapter{Always, Sometimes, Never, Now, Later}
\label{ch:ru-always-sometimes-never-now-later}
\hlchaptermodality{74}

\begin{chapteropening}
\textbf{By the end of this chapter:} \emph{I can say always, sometimes, never, now and later.}

Complete the last lesson of chapter 74: I can say always, sometimes, never, now and later.
\end{chapteropening}

\section[\texorpdfstring{\ru{всегда} (vsegdá)}{всегда (vsegdá)}]{\ru{всегда} (vsegdá) — always}

\label{lesson:RU-C74-always}



\begin{quote}
Before the new one: say the Russian for a lot, then the Russian for a little.
\end{quote}

\subsection*{You'll want to know: \ru{всегда}}
\textbf{\ru{всегда}} (\emph{vsegdá}) — \textquotedblleft{}always\textquotedblright{}.

\begin{cousinweb}
\textbf{\ru{всегда}} is \emph{\ru{весь}}, \textquotedblleft{}all\textquotedblright{}, with \emph{-\ru{гда}}, the time ending of \emph{\ru{когда}}, \textquotedblleft{}when\textquotedblright{}, and \emph{\ru{тогда}}, \textquotedblleft{}then\textquotedblright{}. Lithuanian \emph{visadà}, \textquotedblleft{}always\textquotedblright{}, matches the forms without \emph{\ru{г}} in other Slavic languages, such as Czech \emph{vezdy}; Vasmer ties \emph{-\ru{гда}} itself to Lithuanian \emph{kadà}, \textquotedblleft{}when\textquotedblright{}.
\end{cousinweb}

\begin{grammarlens}[title={What you've built}]
One more word for this chapter.
\end{grammarlens}

\subsection*{Your turn}
\begin{itemize}
\raggedright
  \item \emph{Say it:} \emph{vsegdá}
  \item \emph{Say it:} \emph{vsegdá}, once more
  \item \emph{Say it:} say \emph{vsegdá}
  \item \emph{Recall:} say the Russian for a lot, then the Russian for a little, then say \emph{vsegdá} again
\end{itemize}

\begin{tcolorbox}[breakable,colback=teal!4,colframe=teal!35!black,title={Before you move on}]
What does \ru{всегда} mean? (\textquotedblleft{}Always\textquotedblright{}.) Say it once more.
\end{tcolorbox}
\section[\texorpdfstring{\ru{иногда} (inogdá)}{иногда (inogdá)}]{\ru{иногда} (inogdá) — sometimes}

\label{lesson:RU-C74-sometimes}



\begin{quote}
Before the new one: say the Russian for a little, then the Russian for always.
\end{quote}

\subsection*{You'll want to know: \ru{иногда}}
\textbf{\ru{иногда}} (\emph{inogdá}) — \textquotedblleft{}sometimes\textquotedblright{}.

\begin{cousinweb}
\textbf{\ru{иногда}} is \emph{\ru{иной}}, \textquotedblleft{}other\textquotedblright{}, with the time ending \emph{-\ru{гда}} of \emph{\ru{когда}}, \textquotedblleft{}when\textquotedblright{} (Vasmer); Old Russian \emph{\ru{иногда}} meant \textquotedblleft{}once, at another time\textquotedblright{}.
\end{cousinweb}

\begin{grammarlens}[title={What you've built}]
One more word for this chapter.
\end{grammarlens}

\subsection*{Your turn}
\begin{itemize}
\raggedright
  \item \emph{Say it:} \emph{inogdá}
  \item \emph{Say it:} \emph{inogdá}, once more
  \item \emph{Say it:} say \emph{inogdá}
  \item \emph{Recall:} say the Russian for a little, then the Russian for always, then say \emph{inogdá} again
\end{itemize}

\begin{tcolorbox}[breakable,colback=teal!4,colframe=teal!35!black,title={Before you move on}]
What does \ru{иногда} mean? (\textquotedblleft{}Sometimes\textquotedblright{}.) Say it once more.
\end{tcolorbox}
\section[\texorpdfstring{\ru{никогда} (nikogdá)}{никогда (nikogdá)}]{\ru{никогда} (nikogdá) — never}

\label{lesson:RU-C74-never}



\begin{quote}
Before the new one: say the Russian for always, then the Russian for sometimes.
\end{quote}

\subsection*{You'll want to know: \ru{никогда}}
\textbf{\ru{никогда}} (\emph{nikogdá}) — \textquotedblleft{}never\textquotedblright{}. Built on \textbf{\ru{когда}}, when.

\begin{cousinweb}
\textbf{\ru{никогда}} is \emph{\ru{ни}}, \textquotedblleft{}not even\textquotedblright{}, with \emph{\ru{когда}}, \textquotedblleft{}when\textquotedblright{}: \textquotedblleft{}at no time\textquotedblright{}. Its \emph{\ru{к}-} is the old question stem of \emph{\ru{кто}}, \textquotedblleft{}who\textquotedblright{}.
\end{cousinweb}

\begin{grammarlens}[title={What you've built}]
One more word for this chapter.
\end{grammarlens}

\subsection*{Your turn}
\begin{itemize}
\raggedright
  \item \emph{Say it:} \emph{nikogdá}
  \item \emph{Say it:} \emph{nikogdá}, once more
  \item \emph{Say it:} say \emph{nikogdá}
  \item \emph{Recall:} say the Russian for always, then the Russian for sometimes, then say \emph{nikogdá} again
\end{itemize}

\begin{tcolorbox}[breakable,colback=teal!4,colframe=teal!35!black,title={Before you move on}]
What does \ru{никогда} mean? (\textquotedblleft{}Never\textquotedblright{}.) Say it once more.
\end{tcolorbox}
\section[\texorpdfstring{\ru{сейчас} (seychás)}{сейчас (seychás)}]{\ru{сейчас} (seychás) — now, right now}

\label{lesson:RU-C74-now}



\begin{quote}
Before the new one: say the Russian for sometimes, then the Russian for never.
\end{quote}

\subsection*{You'll want to know: \ru{сейчас}}
\textbf{\ru{сейчас}} (\emph{seychás}) — \textquotedblleft{}now, right now\textquotedblright{}.

\begin{cousinweb}
\textbf{\ru{сейчас}} is \emph{\ru{сей}}, an old word for \textquotedblleft{}this\textquotedblright{}, with \emph{\ru{час}}, \textquotedblleft{}hour\textquotedblright{}: \textquotedblleft{}this hour\textquotedblright{}.
\end{cousinweb}

\begin{grammarlens}[title={What you've built}]
One more word for this chapter.
\end{grammarlens}

\subsection*{Your turn}
\begin{itemize}
\raggedright
  \item \emph{Say it:} \emph{seychás}
  \item \emph{Say it:} \emph{seychás}, once more
  \item \emph{Say it:} say \emph{seychás}
  \item \emph{Recall:} say the Russian for sometimes, then the Russian for never, then say \emph{seychás} again
\end{itemize}

\begin{tcolorbox}[breakable,colback=teal!4,colframe=teal!35!black,title={Before you move on}]
What does \ru{сейчас} mean? (\textquotedblleft{}Now, right now\textquotedblright{}.) Say it once more.
\end{tcolorbox}
\section[\texorpdfstring{\ru{потом} (potóm)}{потом (potóm)}]{\ru{потом} (potóm) — later, then}

\label{lesson:RU-C74-later}



\begin{quote}
Before the new one: say the Russian for never, then the Russian for now.
\end{quote}

\subsection*{You'll want to know: \ru{потом}}
\textbf{\ru{потом}} (\emph{potóm}) — \textquotedblleft{}later, then\textquotedblright{}.

\begin{cousinweb}
\textbf{\ru{потом}} is \emph{\ru{по}}, \textquotedblleft{}after\textquotedblright{}, with \emph{\ru{том}}, a form of \emph{\ru{тот}}, \textquotedblleft{}that\textquotedblright{}: \textquotedblleft{}after that\textquotedblright{}.
\end{cousinweb}

\begin{grammarlens}[title={What you've built}]
One more word for this chapter.
\end{grammarlens}

\subsection*{Your turn}
\begin{itemize}
\raggedright
  \item \emph{Say it:} \emph{potóm}
  \item \emph{Say it:} \emph{potóm}, once more
  \item \emph{Say it:} say \emph{potóm}
  \item \emph{Recall:} say the Russian for never, then the Russian for now, then say \emph{potóm} again
\end{itemize}

\begin{tcolorbox}[breakable,colback=teal!4,colframe=teal!35!black,title={Before you move on}]
What does \ru{потом} mean? (\textquotedblleft{}Later, then\textquotedblright{}.) Say it once more.
\end{tcolorbox}
